% !TEX program = xelatex
% Отчёт о текущем состоянии работы.
% Сборка:  latexmk -xelatex report.tex   (или дважды xelatex report.tex)
\documentclass[11pt,a4paper]{article}

\usepackage{fontspec}
\usepackage{polyglossia}
\setmainlanguage{russian}
\setotherlanguage{english}

% Шрифты с поддержкой кириллицы (есть в стандартной поставке Windows и TeX Live).
\setmainfont{DejaVu Serif}
\setsansfont{DejaVu Sans}
\setmonofont{DejaVu Sans Mono}

\usepackage{geometry}
\geometry{margin=2.3cm}

\usepackage{amsmath}
\usepackage{booktabs}
\usepackage{array}
\usepackage{enumitem}
\usepackage{xcolor}
\usepackage{caption}
\usepackage{microtype}
\usepackage{fvextra}
\usepackage{titlesec}
\usepackage[hidelinks]{hyperref}

% Немного запасной растяжимости, чтобы длинные идентификаторы в \code
% реже вылезали за правое поле.
\setlength{\emergencystretch}{2em}

% Заголовки: раньше в оглавлении номер из двух цифр почти слипался с текстом.
\titleformat{\section}{\normalfont\Large\bfseries}{\thesection}{0.8em}{}
\titleformat{\subsection}{\normalfont\large\bfseries}{\thesubsection}{0.7em}{}

\captionsetup{font=small,labelfont=bf}

% Номер раздела из двух цифр («10») не должен слипаться с заголовком
% в оглавлении: расширяем поле под номер.
\makeatletter
\renewcommand*\l@section[2]{%
  \ifnum \c@tocdepth >\z@
    \addpenalty\@secpenalty
    \addvspace{1.0em \@plus\p@}%
    \setlength\@tempdima{2.6em}%
    \begingroup
      \parindent \z@ \rightskip \@pnumwidth
      \parfillskip -\@pnumwidth
      \leavevmode \bfseries
      \advance\leftskip\@tempdima
      \hskip -\leftskip
      #1\nobreak\hfil\nobreak\hbox to\@pnumwidth{\hss #2}\par
    \endgroup
  \fi}
\makeatother

% --- листинги ---------------------------------------------------------------
% Используется fvextra (расширение fancyvrb): в XeLaTeX пакет listings
% переставляет открывающую скобку после кириллического слова (известный
% дефект), а fancyvrb выводит кириллицу корректно.
\definecolor{codebg}{RGB}{247,247,249}
\definecolor{codeframe}{RGB}{210,210,215}

\fvset{
  fontsize=\small,
  frame=single,
  framesep=4pt,
  rulecolor=\color{codeframe},
  bgcolor=codebg,
  xleftmargin=6pt,
  xrightmargin=4pt,
}

% Команды оболочки (допускают перенос длинных строк).
\DefineVerbatimEnvironment{ShellCode}{Verbatim}{%
  breaklines=true,
  breakanywhere=true,
}
% Дерево каталогов (переносы не нужны).
\DefineVerbatimEnvironment{TreeCode}{Verbatim}{%
  breaklines=false,
}
% Программа на Python.
\DefineVerbatimEnvironment{PyCode}{Verbatim}{%
  breaklines=true,
  breakanywhere=true,
}

\newcommand{\code}[1]{\texttt{#1}}

\title{\bfseries Отчёт о текущем состоянии работы\\[4pt]
\large Реорганизация проекта и первый этап: детектор речевой активности}
\author{Магистерская работа: диагностика свойств интеллекта\\
по временным характеристикам речи}
\date{05.11.2026}

\begin{document}
\sloppy
\maketitle

\begin{abstract}
\noindent
Прежняя версия проекта (\code{voice\_signals}) переработана в более
компактный проект \code{magister\_work} с единой схемой обработки:
сигнал $\rightarrow$ сегментация $\rightarrow$ признак. Реализована программа
первого этапа --- детектор речевой активности (VAD, от англ.\ Voice Activity
Detection): перенесён и упрощён энергетический метод, добавлен собственный
метод на основе энтропии заполнения фазового пространства (метод научного
руководителя), реализованы метрики качества и 22 детерминированных теста.
Все тесты проходят; оба метода проверены на реальном (не целевом) речевом
файле. Основное ограничение сохраняется: целевой набор записей пока не
получен.
\end{abstract}

\tableofcontents
\bigskip

% ============================================================================
\section{Резюме}
% ============================================================================

\begin{itemize}[leftmargin=1.4em]
  \item Проект реорганизован: вместо разросшегося \code{voice\_signals}
        создан компактный \code{magister\_work}. Из прежней программы
        перенесены только рабочие части; вспомогательные средства (проверка
        параметров файлов, извлечение признаков, экспериментальные сценарии)
        не переносились.
  \item Первый этап выполнен: построен детектор речевой активности с двумя
        методами и проверяемым качеством.
  \item Реализованы покадровые метрики, включая основную метрику
        \code{false\_pause\_precision}.
  \item Тесты: все 22 проходят. Зависимости минимальны (\code{pytest},
        \code{PyYAML}); используется Python~3.13.
  \item Ограничение: \textbf{целевой набор записей респондентов не получен},
        поэтому настройка порогов и исследование связи с результатами
        когнитивных тестов отложены.
\end{itemize}

% ============================================================================
\section{Цель работы и общая схема}
% ============================================================================

Научная цель --- проверить и развить гипотезу о связи временной организации
речи с уровнем интеллекта. Опора --- работа Кузенкова, Пака и Французенко
(тесты Айзенка, чтение и пересказ рассказа): с ростом IQ укорачиваются паузы
при чтении и длительности звуков при пересказе, а линейная модель по
временным параметрам восстанавливает усреднённый IQ.

Работа разбита на три этапа (таблица~\ref{tab:stages}).

\begin{table}[h]
\centering
\small
\caption{Этапы работы.}
\label{tab:stages}
\begin{tabular}{@{}c p{4.2cm} p{4.2cm} p{3.2cm}@{}}
\toprule
Этап & Содержание & Результат & Состояние \\
\midrule
\textbf{1} & Сегментация речи и паузы & детектор и метрики & \textbf{выполнен} \\
2 & Временные и нелинейные признаки & вектор признаков на запись & частично \\
3 & Связь признаков с IQ & статистическая модель & ожидает данных \\
\bottomrule
\end{tabular}
\end{table}

Первый этап определён научным руководителем как обязательный и
предварительный: он даёт основу для всех последующих этапов. Именно ему
посвящена основная часть отчёта.

% ============================================================================
\section{Структура проекта}
% ============================================================================

\begin{TreeCode}
magister_work/
  src/speechdyn/            # библиотека
    contracts.py            # SegmentRecord, разметка, разбиение
    framing.py              # TimeGrid, кадры, параболическая огибающая
    wav.py                  # чтение mono 16-bit PCM WAV
    config.py               # YAML и хеш конфигурации
    metrics.py              # точность, полнота, F1, IoU
    dynamics.py             # модифицированный R/S: H и T (варианты A и C)
    synthetic.py            # синтетический сигнал для тестов
    vad/
      base.py               # VadBackend, состояния -> отрезки
      energy.py             # EnergyVad (гистерезис)
      entropy.py            # EntropyVad (энтропия фазового пространства)
  scripts/run_vad.py        # запуск методов на файле
  tests/test_stage1.py      # 22 теста
  configs/default.yaml      # параметры VAD, признаков, разметки
  docs/                     # план, спецификация, идеи, этот отчёт
  wiki/                     # исходные статьи и их извлечения
\end{TreeCode}

Назначение модулей приведено в таблице~\ref{tab:modules} (краткая справка по
источникам --- в разделе~\ref{sec:sources}).

\begin{table}[h]
\centering
\small
\caption{Модули библиотеки \code{speechdyn}.}
\label{tab:modules}
\begin{tabular}{@{}p{4.4cm}p{9.3cm}@{}}
\toprule
Модуль & Назначение \\
\midrule
\code{contracts.py} & Единая структура \code{SegmentRecord}; разметка
  \code{speech}/\code{silent\_pause}/\code{hesitation}/\code{artifact};
  достройка полного разбиения на речь и не-речь. \\
\code{framing.py} & Временная сетка \code{TimeGrid}, кадры,
  параболическая огибающая (метод 2021 г.). \\
\code{wav.py} & Строгое чтение mono 16-bit PCM WAV; иные форматы отклоняются. \\
\code{config.py} & Загрузка YAML и устойчивый хеш конфигурации. \\
\code{metrics.py} & Покадровые метрики и совпадение отрезков. \\
\code{dynamics.py} & Модифицированный R/S: показатель Хёрста $H$ и
  характерное время $T$. \\
\code{synthetic.py} & Детерминированный синтетический сигнал с эталонной
  разметкой. \\
\code{vad/base.py} & Абстрактный класс \code{VadBackend} и общая
  постобработка. \\
\code{vad/energy.py} & Энергетический метод с гистерезисом. \\
\code{vad/entropy.py} & Метод на основе энтропии заполнения фазового
  пространства. \\
\bottomrule
\end{tabular}
\end{table}

% ============================================================================
\section{Первый этап: детектор речевой активности}
% ============================================================================

Задача этапа: по монофонической записи в формате WAV построить
непересекающиеся отрезки речи и представить их в единой структуре
\code{SegmentRecord}. Классификация бинарная
(\code{speech}/\code{non\_speech}); полная разметка из четырёх меток
используется при ручном анализе и на последующих этапах.

Формулировка результата намеренно не когнитивная: \emph{программа размечает
речевые события}, а признаки затем могут исследоваться на связь с
результатами когнитивных тестов.

\subsection{Энергетический метод (базовый)}

Порядок расчёта: среднеквадратичное значение сигнала в кадрах длиной 20~мс с
шагом 10~мс; переход к децибелам; порог шума как 20-й процентиль; гистерезис
(порог включения 10~дБ, порог выключения 6~дБ); постобработка (минимальная
длительность речи 80~мс, объединение пауз короче 80~мс). Метод
детерминирован и не требует внешних библиотек.

\subsection{Метод на основе энтропии фазового пространства}

Реализация метода из сборника «Нейроинформатика-2021»:

\begin{enumerate}[leftmargin=1.6em]
  \item \textbf{огибающая} $E$: взвешивание отсчётов окна длиной~10~мс
        параболой (нуль на краях, единица в центре) и выбор максимума;
  \item \textbf{энтропия фазового пространства}: траектория
        $\{(x_n, x_{n+k})\}$, фазовую плоскость делим на $N\times N$
        квадратов, вероятности $p_i = m_i / M_k$, энтропия
        $S_k = -\sum_i p_i \log_2 p_i$; задержки $k = 2, \dots, 16$;
  \item \textbf{признаки окна} (длительность $100$~мс, шаг $25$~мс):
        $(\ln E,\ \langle S \rangle,\ \sigma(S))$;
  \item \textbf{решающее правило}: окно считается паузой, если
        \[
          \ln E < -5.0
          \quad\text{или}\quad
          (3.43 - \langle S\rangle)^2 + (0.0653 - \sigma(S))^2 < 0.0288 ;
        \]
  \item \textbf{постобработка} --- как в энергетическом методе.
\end{enumerate}

Пороги перенесены в конфигурацию как \emph{начальные} (они получены авторами
для записей с отношением сигнал/шум не менее 35~дБ) и подлежат настройке на
целевом наборе записей.

\subsection{Проверка на реальной речи}

Оба метода запущены на реальном (не целевом) файле \code{Ru-Vegetables.wav}
(частота дискретизации 16~кГц, длительность 124.5~с) без подбора параметров;
результаты приведены в таблице~\ref{tab:realcheck}.

\begin{table}[h]
\centering
\small
\caption{Результат на реальном файле \code{Ru-Vegetables.wav}.}
\label{tab:realcheck}
\begin{tabular}{@{}lrr@{}}
\toprule
Метод & Число отрезков & Доля речи \\
\midrule
энергетический & 79 & 70.8\,\% \\
энтропийный    & 82 & 72.8\,\% \\
\bottomrule
\end{tabular}
\end{table}

Близкие значения говорят лишь о разумной согласованности двух независимых
подходов на начальном уровне. Это \emph{не} оценка качества: эталонной
разметки для этого файла нет.

% ============================================================================
\section{Метрики качества}
% ============================================================================

Для оценки на общей временной сетке зафиксируем две двоичные
последовательности: $p_i$ --- предсказание программы, $r_i$ --- эталон
($1$ --- речь, $0$ --- пауза). Сравнение по кадрам даёт числа
истинно-положительных ($TP$), ложноположительных ($FP$) и
ложноотрицательных ($FN$) кадров. Тогда

\[
  P = \frac{TP}{TP+FP}, \qquad
  R = \frac{TP}{TP+FN}, \qquad
  F_1 = \frac{2PR}{P+R}, \qquad
  \mathrm{IoU} = \frac{TP}{TP+FP+FN},
\]

где $P$ --- точность, $R$ --- полнота, $F_1$ --- $F_1$-мера,
$\mathrm{IoU}$ --- коэффициент пересечения (от англ.\ Intersection over
Union). Основная метрика \code{false\_pause\_precision} считается отдельно
по паузам: это доля верных пауз среди всех пауз, предсказанных программой
(отношение истинно-отрицательных кадров к их сумме с ложноотрицательными).
Эталон приводится к сетке функцией \code{rasterize}. Полный набор метрик
приведён в таблице~\ref{tab:metrics}.

\begin{table}[h]
\centering
\small
\caption{Метрики качества.}
\label{tab:metrics}
\begin{tabular}{@{}lll@{}}
\toprule
Метрика & Смысл & Роль \\
\midrule
\code{false\_pause\_precision} & доля верных пауз среди предсказанных & \textbf{основная} \\
\code{speech\_recall}          & полнота речи & контроль \\
\code{pause\_duration\_bias\_sec} & смещение суммарной длительности речи & контроль \\
\code{precision}/\code{recall}/\code{f1}/\code{iou} & стандартные покадровые метрики & контроль \\
\code{segment\_f1}             & совпадение отрезков с допуском $\pm100$~мс & контроль \\
\bottomrule
\end{tabular}
\end{table}

Метрики, полученные на синтетическом сигнале, характеризуют только
работоспособность программы, но не качество на реальных записях.

% ============================================================================
\section{Признаки временной динамики}
% ============================================================================

Перенесён модифицированный R/S-анализ (Кузенков, Логинов, 2014): показатель
Хёрста $H$ и характерное время $T$. Сохранены два явно разделённых способа
расчёта:

\begin{itemize}[leftmargin=1.4em]
  \item \code{full\_record} (вариант~A) --- огибающая всей записи вместе с
        паузами, окно $100$~мс, $\tau$ до $N$; значения сопоставимы со
        статьёй;
  \item \code{segment} (вариант~C) --- на отдельном речевом отрезке;
        вспомогательный признак, помеченный как несопоставимый между
        отрезками разной длины и со статьёй.
\end{itemize}

Несмешиваемость закреплена в самой структуре результата (поле \code{scope} и
признак сопоставимости), а не только в документации. Временные признаки
второго этапа (ряды длительностей звуков $V$ и пауз $U$, их статистики,
скорость речи) пока не реализованы --- они подключаются поверх результата
сегментации.

% ============================================================================
\section{Тесты и воспроизводимость}
% ============================================================================

Набор из 22 детерминированных тестов покрывает следующее
(таблица~\ref{tab:tests}).

\begin{table}[h]
\centering
\small
\caption{Проверяемые свойства.}
\label{tab:tests}
\begin{tabular}{@{}p{4.2cm}p{9.3cm}@{}}
\toprule
Группа & Что проверяется \\
\midrule
структуры данных & границы отрезков, метки, пересечения, достройка разбиения \\
кадры и огибающая & временная сетка, среднеквадратичные значения, огибающая \\
состояния $\to$ отрезки & отсев короткой речи, объединение короткой паузы \\
энергетический метод & находит речь; тишина даёт пустой результат; времена
  лежат в пределах файла \\
энтропийный метод & энтропия константы равна нулю; диапазон значений; форма
  признаков \\
метрики & полное совпадение; \code{false\_pause\_precision}; приведение к
  сетке \\
R/S & признаки \code{full\_record}/\code{segment}; короткий вход недопустим \\
\bottomrule
\end{tabular}
\end{table}

Воспроизводимость: генератор синтетического сигнала детерминирован
(фиксированное начальное состояние), все расчёты детерминированы, а
конфигурация имеет устойчивый хеш, который попадает в отчёты.

% ============================================================================
\section{Используемые источники}
\label{sec:sources}
% ============================================================================

Исходные статьи собраны в каталоге \code{wiki/} (файлы PDF и их текстовые
извлечения). Каждый источник даёт свой круг идей (таблица~\ref{tab:sources}).

\begin{table}[h]
\centering
\small
\caption{Источники и что берётся из них.}
\label{tab:sources}
\begin{tabular}{@{}p{3.9cm}p{3.1cm}p{6.6cm}@{}}
\toprule
Источник & Авторы, год & Что берётся в работу \\
\midrule
Диагностика свойств интеллекта по временным характеристикам речи &
Кузенков, Пак, Французенко &
Ряды длительностей звуков $V$ и пауз $U$, их медиана, среднее и стандартное
отклонение, скорость речи; нормировка, метод главных компонент, линейная
модель оценки IQ. Основа научной идеи (этапы~2--3). \\[2pt]
Использование метода нормированного размаха \dots &
Кузенков, Логинов, 2014 &
Модифицированный R/S: геометрическая шкала $\tau$, усреднение по
перекрывающимся окнам, показатель Хёрста $H$ и характерное время $T$,
кусочно-линейная аппроксимация, перестановочный тест (этап~2). \\[2pt]
Нейроинформатика, её приложения и анализ данных (2021) &
Кузенков, Логинов &
Собственный детектор речевой активности на основе энтропии заполнения
фазового пространства; параболическая огибающая; решающее правило по
$(\ln E, \langle S\rangle, \sigma(S))$ (этап~1). \\
\bottomrule
\end{tabular}
\end{table}

Отдельно отметим ограничения переноса. Параметры записи (частота
дискретизации 8~кГц, микрофон) и клинические выводы статей не переносятся
автоматически. Показатель $H$ в работе Кузенкова и Логинова близок к единице
($H \approx 0.984 \dots 0.991$ во всех группах) и не различает группы;
различает характерное время $T$.

Полные тексты --- в \code{wiki/docs/}. Краткий разбор идей --- в
\code{docs/02-ideas.md}.

% ============================================================================
\section{Порядок работы с программой}
% ============================================================================

Все команды выполняются из каталога \code{magister\_work}.

\subsection{Установка зависимостей}

\begin{ShellCode}
python -m pip install -r requirements.txt
\end{ShellCode}

Зависимости минимальны: \code{pytest} (тесты) и \code{PyYAML}
(конфигурация). Внешние библиотеки для обработки звука и машинного обучения
на этом этапе не требуются.

\subsection{Запуск тестов}

\begin{ShellCode}
python -m pytest
# Ожидается: 22 passed
\end{ShellCode}

\subsection{Запуск детектора на файле}

\begin{ShellCode}
# оба метода (энергетический и энтропийный)
python scripts/run_vad.py path/to/file.wav

# только один метод
python scripts/run_vad.py path/to/file.wav --backend energy
python scripts/run_vad.py path/to/file.wav --backend entropy
\end{ShellCode}

Программа выводит число отрезков, долю времени речи и границы отрезков.
Входной файл должен быть монофоническим WAV с 16-битными отсчётами; при
другом формате выводится сообщение об ошибке.

\subsection{Программный интерфейс}

\begin{PyCode}
import sys; sys.path.insert(0, "src")
from speechdyn import EnergyVad, EntropyVad, read_wav_mono

audio, sr = read_wav_mono("file.wav")

segments = EnergyVad().process_array(audio, sr, file_id="file")
for s in segments:
    print(f"{s.start_sec:8.3f} {s.end_sec:8.3f} {s.label}")
\end{PyCode}

Для оценки по эталонной разметке служат функции \code{rasterize} и
\code{frame\_metrics} из \code{speechdyn.metrics}; для достройки полного
разбиения на речь и не-речь --- \code{complement\_segments}.

\subsection{Пример запуска на реальном файле}

\begin{ShellCode}
python scripts/run_vad.py ..\voice_signals\tests\test_data\Ru-Vegetables.wav
\end{ShellCode}

\subsection{Сборка отчёта}

\begin{ShellCode}
cd docs
latexmk -xelatex report.tex
# результат: report.pdf
\end{ShellCode}

% ============================================================================
\section{Ограничения и допущения}
% ============================================================================

\begin{itemize}[leftmargin=1.4em]
  \item \textbf{Целевой набор записей не получен} (дата доступа неизвестна).
        Следствие: пороги детектора не настроены, а исследование связи
        признаков с IQ отложено. На промежуточном этапе используются образцы
        реальной (не целевой) речи для разведочного анализа.
  \item Разметку выполняет \textbf{один человек} (научный руководитель);
        объём --- десятки файлов, поэтому повторная разметка и оценка
        согласованности не применяются.
  \item Правовых ограничений нет; остаётся ограничение по средствам ---
        только локальный запуск и бесплатный тариф облачных вычислений.
  \item В программе явно помечены допущения:
\end{itemize}

\begin{ShellCode}
ASSUMPTION: sample_rate 16000 -- до аудита целевого набора записей.
ASSUMPTION: пороги энтропийного метода (-5.0; 3.43/0.0653; 0.0288)
            перенесены из статьи при SNR >= 35 дБ и не настроены.
UNKNOWN:     фактические формат, число каналов, частота дискретизации,
             длительности записей.
TODO_AFTER_DATA_AUDIT: зафиксировать параметры и настроить пороги.
\end{ShellCode}

% ============================================================================
\section{Ближайшие задачи}
% ============================================================================

\begin{enumerate}[leftmargin=1.6em]
  \item Добавить метод сравнения WebRTC (интерфейс \code{VadBackend} готов).
  \item Подготовить инструкцию по разметке для одного человека (разметка уже
        согласована).
  \item Проверить устойчивость к шуму, ограничению сигнала и разной громкости.
  \item Этап~2: реализовать временные признаки (ряды $V$ и $U$, скорость
        речи) поверх результата сегментации.
  \item Составить разведочный отчёт по метрикам на реальных (не целевых)
        образцах.
\end{enumerate}

Открытые вопросы к научному руководителю: порядок включения нелинейных
признаков; достаточно ли кусочно-линейной аппроксимации для времени $T$ или
нужна кусочно-полиномиальная; какой набор когнитивных тестов использовать и
на каком этапе проверять связь с IQ.

\bigskip
\hrule
\bigskip
\noindent\footnotesize
Проект: \code{magister\_work}, библиотека \code{speechdyn}. Среда: Python~3.13,
XeLaTeX (TeX~Live~2026), шрифты DejaVu. Дата: 05.11.2026.

\end{document}
